\documentclass[12pt]{article}
\usepackage[utf8]{inputenc}
\usepackage[T1]{fontenc}
\usepackage{hyperref}
\title{Yangum Gel Language Resources for K-12 Curriculum Development}
\author{Compiled by Claude Sonnet 5.5 (claude-sonnet-5-5) for Andr\'as Kornai}
\date{}
\begin{document}
\maketitle

\section{Introduction}
Yangum Gel (ISO 639-3 \texttt{ygl}) is a Torricelli language of Papua New Guinea. It belongs to the Yangum branch, alongside Yangum Dey and Yangum Mon. Wikipedia calls it nearly extinct. ASJP lists about 45 speakers, while the 2000 census figure of 960 on Wikipedia covers the whole Yangum cluster. Sources disagree on the glottocode: the ASJP entry for Yangum Gel is cross-referenced to \texttt{yang1295}, and Wikipedia cites \texttt{yang1303} for the broader cluster. Little has been published, and I found no teaching materials.

\section{Primary Educational Resources}
None found.

\section{Linguistic and Reference Materials}
\begin{itemize}
\item Wikipedia, Yangum language (covers the cluster): \url{https://en.wikipedia.org/wiki/Yangum_language}
\item ASJP wordlist, Yangum Gel: \url{https://asjp.clld.org/languages/YANGUM_GEL}
\item Crossgram language page (\texttt{yang1295}): \url{https://crossgram.clld.org/languages/yang1295}
\item Glottolog entry: \url{https://glottolog.org/resource/languoid/id/heyo1241}
\item Wikipedia, Torricelli languages (family context): \url{https://en.wikipedia.org/wiki/Torricelli_languages}
\end{itemize}

\section{Religious and Community Resources}
None found for Yangum Gel specifically.

\section{Digital and Audio Resources}
None found.

\section{Curriculum Development Context}
The language is severely endangered, with few speakers by the ASJP count, and no school use is known. Any curriculum would start with community-led documentation: word lists, recordings and a practical orthography. It would draw on the Yangum cluster and the wider Torricelli literature, and should wait for fresh fieldwork. No ready-made native-language materials were found.

\end{document}